Modern neural networks suffer from over-parameterization and require substantial compute and memory during both training and inference. While many pruning techniques exist, most rely on post-hoc weight magnitude heuristics or static pre-training sensitivity metrics. In this paper, we propose \textbf{Dynamic Activity-Dependent Pruning (DADP)}, a biologically inspired structural plasticity mechanism. DADP evaluates connection importance during training using the expected product of pre-synaptic activation and post-synaptic loss gradient: $I_{ij} = \mathbb{E}\left[ \left| a_i \cdot \frac{\partial L}{\partial y_j} \right| \right]$. This creates a self-regulating feedback loop where the network dynamically prunes redundant connections and allows surviving pathways to recover. Importantly, DADP does not require pre-set target sparsities; instead, the network organically finds its own sparsity equilibrium while naturally inducing channel and neuron-level pruning. Across MLP, VGG-16, ResNet-18, BiLSTM-CRF, and MiniBERT architectures on MNIST, CIFAR-10, CoNLL-2003, and SST-2, DADP consistently matches or outperforms state-of-the-art sparse training methods (such as Magnitude pruning, SNIP, and RigL), retaining up to 79.91\% accuracy at 99\% sparsity on ResNet-18. Furthermore, we analyze representation quality using Matrix-Based Shannon Entropy and Effective Rank, showing that DADP preserves layerwise information even under extreme compression. DADP provides a simple, single-pass training framework that achieves high model compression and strong generalization without expensive prune-and-retrain cycles.
